% !TEX root = ../Thesis.tex
\chapter{Introduction}
\label{ch:introduction}

\begingroup
\setaftersecskip{6pt}
\setlength{\emergencystretch}{2em}
Digital collections are often shown on websites or in lists. These formats make direct searches easy, but they do not feel like an exhibition. This thesis studies how retrieved media can instead be presented and explored in a virtual museum that users can move through freely.

\section{Motivation}
\label{sec:introduction-motivation}

A digital collection is accessible from any place and at any time. Standard search interfaces focus on an overview and direct access to individual objects. They usually present media next to each other in a two-dimensional format. As a result, there is little sense of space or connection between the objects. A virtual museum can present the same objects in a spatial environment while keeping the advantages of digital access.

Since physical museums have limited exhibition space, they can only show a small part of a large collection at once. Virtual museums do not have this limitation and allow a much larger number of objects to be accessed without requiring the original items to be displayed. The objects can remain safely stored, while their digital representations do not age or decay. The same representation can appear in several exhibitions at once, and one virtual exhibition can even bring together objects from different institutions, although the originals are stored in completely different locations.

Virtual Reality extends this concept by allowing visitors to experience search results as exhibits. Visitors can walk around a 3D model, inspect it from different angles, and use it as the starting point for a new search. In this way, each exhibition becomes a room to explore, and an exhibit can determine where the exploration goes next.

\enlargethispage{\baselineskip}
\section{Initial Situation and Problem Statement}
\label{sec:introduction-problem}

Existing systems already show that virtual environments can support museum visits and the presentation of digital objects. Most of them, however, focus on a specific search method, media type, or interaction. This thesis shows how text-based and similarity-based search can be combined with the presentation of images, videos, and 3D models in a single VR exhibition.

This combination introduces two challenges. First, images, videos, and 3D models have different characteristics. They cannot use the same form of presentation, even though they should all belong to one exhibition. Second, the hardware of the tested platforms varies. The application can run on the headset itself or be streamed from a PC. Users can also interact with controllers or hand tracking. Even with these differences, the process in the museum should stay the same.

This leads to the following research question:

\begin{quote}
\emph{How can a multimodal collection be shown as a dynamic VR exhibition using vitrivr-engine and explored through text and similarity searches?}
\end{quote}

Answering this question requires both a concept and a prototype. The concept describes how search results become an exhibition and how a visitor can continue the search from an exhibited object. To implement this concept, the VR application connects to the existing search backend. It loads results and shows each media object in a suitable form. Another search can then be started from an exhibited object. These steps should work in the same way on all tested devices.

\section{Project Idea}
\label{sec:introduction-project-idea}

The project uses one virtual museum room whose exhibition changes with every search. Images and videos returned by a search are displayed on the walls, while 3D models are placed in areas on the floor. Running another query replaces the exhibits without rebuilding the room itself.

Every exhibit has a \emph{Similar} action. When a visitor selects it, the selected exhibit starts a new similarity search. The new results are then shown as the next exhibition. This allows visitors to explore the collection step by step. The room is not based on a real museum. Instead, it is a spatial interface for the digital collection.

\section{Scope and Contributions}
\label{sec:introduction-scope}

\paragraph{Scope} This thesis studies a VR prototype with a prepared collection of images, videos, and 3D models. It offers text and similarity searches and can be controlled with controllers or hand tracking. On the Meta Quest~3 and VIVE Focus~3, it runs directly as a standalone application. A Windows version can also be streamed to a supported headset. The work also includes the connection to vitrivr, installation instructions, and a usability evaluation with a small group of participants. New search or embedding models are not part of this thesis. The same applies to a production-ready system, support for all VR devices, and a scan of a real museum. The prototype uses existing components such as the vitrivr-engine, the Descriptor Server, CLIP, PostgreSQL, and pgvector.

\pagebreak
\paragraph{Contributions} The contributions of this thesis are
\begin{enumerate*}[label=(\roman*), itemjoin={{, }}, itemjoin*={{, and }}]
  \item a concept that transfers a ranked result set into a spatial exhibition
  \item a VR prototype that shows images, videos, and 3D models together and can be controlled with controllers or hand tracking
  \item a connection between the prototype and the vitrivr-engine, including the presentation of the media returned by a search
  \item a usability evaluation with a small group of participants, together with instructions for setting up the complete system.
\end{enumerate*}

So the main contributions are the VR interface and the interaction between search, media, and exhibition.

\section{Thesis Outline}
\label{sec:introduction-outline}

The thesis begins with the background in Chapter~\ref{ch:background}. It introduces Extended Reality, multimedia retrieval, vector representations, CLIP, 3D model search, and the vitrivr-engine. Chapter~\ref{ch:related-work} looks at related work on virtual exhibitions and immersive search systems. Chapter~\ref{ch:concepts} introduces the design of the VR museum, and Chapter~\ref{ch:implementation} explains how it was implemented. Chapter~\ref{ch:evaluation} reports the evaluation and its results. Chapter~\ref{ch:conclusion} closes the thesis by answering the research question, discussing the work, and suggesting future improvements.

\endgroup
